% 综合性能对比柱状图
\begin{tikzpicture}[scale=0.85, transform shape]
  % 参数定义
  \def\barwidth{0.55}
  \def\gap{0.25}
  \def\groupwidth{2.2}

  % ===== 坐标轴 =====
  \draw[-Stealth, thick] (-0.5, 0) -- (12.5, 0);
  \draw[-Stealth, thick] (0, -0.3) -- (0, 7.5) node[above, font=\footnotesize] {指标值};

  % Y轴刻度
  \foreach \y/\lbl in {0/0, 1.5/25, 3/50, 4.5/75, 6/100} {
    \draw[gray!50, thin] (0.1, \y) -- (12, \y);
    \draw[thick] (-0.1, \y) -- (0.1, \y) node[left, font=\tiny] {\lbl};
  }

  % ===== 指标1：规划成功率（%） =====
  \def\xbase{1.2}
  % Apollo原始: 79.8%
  \fill[dangerred!60, draw=dangerred!80, thick] (\xbase-\barwidth/2, 0) rectangle (\xbase+\barwidth/2, 4.788);
  \node[font=\tiny, dangerred!80!black] at (\xbase, 5.05) {79.8};
  % 改进: 95.6%
  \fill[lightgreen!70, draw=lightgreen!80!black, thick] (\xbase+\gap+\barwidth/2, 0) rectangle (\xbase+\gap+3*\barwidth/2, 5.736);
  \node[font=\tiny, lightgreen!70!black] at (\xbase+\gap+\barwidth, 6.0) {95.6};
  % 标签
  \node[font=\tiny, align=center] at (\xbase+\gap/2+\barwidth/2, -0.6) {规划成功率\\(\%)};

  % ===== 指标2：通行速度（m/s） =====
  \def\xbase{3.6}
  \fill[dangerred!60, draw=dangerred!80, thick] (\xbase-\barwidth/2, 0) rectangle (\xbase+\barwidth/2, 2.865);
  \node[font=\tiny, dangerred!80!black] at (\xbase, 3.13) {3.82};
  \fill[lightgreen!70, draw=lightgreen!80!black, thick] (\xbase+\gap+\barwidth/2, 0) rectangle (\xbase+\gap+3*\barwidth/2, 3.773);
  \node[font=\tiny, lightgreen!70!black] at (\xbase+\gap+\barwidth, 4.04) {5.03};
  \node[font=\tiny, align=center] at (\xbase+\gap/2+\barwidth/2, -0.6) {通行速度\\(m/s)};

  % ===== 指标3：停车次数 =====
  \def\xbase{6.0}
  \fill[dangerred!60, draw=dangerred!80, thick] (\xbase-\barwidth/2, 0) rectangle (\xbase+\barwidth/2, 3.36);
  \node[font=\tiny, dangerred!80!black] at (\xbase, 3.63) {1.68};
  \fill[lightgreen!70, draw=lightgreen!80!black, thick] (\xbase+\gap+\barwidth/2, 0) rectangle (\xbase+\gap+3*\barwidth/2, 0.62);
  \node[font=\tiny, lightgreen!70!black] at (\xbase+\gap+\barwidth, 0.89) {0.31};
  \node[font=\tiny, align=center] at (\xbase+\gap/2+\barwidth/2, -0.6) {停车次数\\(次)};

  % 下降箭头标注
  \draw[-Stealth, thick, deeppink] (\xbase+\gap/2+\barwidth/2, 3.6) -- (\xbase+\gap/2+\barwidth/2, 1.1);
  \node[deeppink, font=\tiny\bfseries] at (\xbase+\gap/2+\barwidth/2+0.6, 2.35) {-81.5\%};

  % ===== 指标4：规划耗时（ms） =====
  \def\xbase{8.4}
  \fill[dangerred!60, draw=dangerred!80, thick] (\xbase-\barwidth/2, 0) rectangle (\xbase+\barwidth/2, 3.12);
  \node[font=\tiny, dangerred!80!black] at (\xbase, 3.39) {26.0};
  \fill[lightgreen!70, draw=lightgreen!80!black, thick] (\xbase+\gap+\barwidth/2, 0) rectangle (\xbase+\gap+3*\barwidth/2, 2.592);
  \node[font=\tiny, lightgreen!70!black] at (\xbase+\gap+\barwidth, 2.86) {21.6};
  \node[font=\tiny, align=center] at (\xbase+\gap/2+\barwidth/2, -0.6) {规划耗时\\(ms)};

  % ===== 指标5：加速度方差 =====
  \def\xbase{10.8}
  \fill[dangerred!60, draw=dangerred!80, thick] (\xbase-\barwidth/2, 0) rectangle (\xbase+\barwidth/2, 3.7);
  \node[font=\tiny, dangerred!80!black] at (\xbase, 3.97) {0.74};
  \fill[lightgreen!70, draw=lightgreen!80!black, thick] (\xbase+\gap+\barwidth/2, 0) rectangle (\xbase+\gap+3*\barwidth/2, 2.35);
  \node[font=\tiny, lightgreen!70!black] at (\xbase+\gap+\barwidth, 2.62) {0.47};
  \node[font=\tiny, align=center] at (\xbase+\gap/2+\barwidth/2, -0.6) {加速度方差\\($\text{m}^2/\text{s}^4$)};

  % ===== 图例 =====
  \fill[dangerred!60, draw=dangerred!80, thick] (3.5, 7.0) rectangle (4.3, 7.35);
  \node[font=\footnotesize, right] at (4.4, 7.175) {Apollo原始};
  \fill[lightgreen!70, draw=lightgreen!80!black, thick] (7.0, 7.0) rectangle (7.8, 7.35);
  \node[font=\footnotesize, right] at (7.9, 7.175) {改进算法};

  % ===== 改进百分比标注（顶部） =====
  \node[deepblue, font=\tiny\bfseries] at (1.6, 6.3) {+15.8\%};
  \node[deepblue, font=\tiny\bfseries] at (4.0, 4.4) {+31.7\%};
  \node[deepblue, font=\tiny\bfseries] at (8.8, 3.7) {-16.9\%};
  \node[deepblue, font=\tiny\bfseries] at (11.2, 4.3) {-36.5\%};
\end{tikzpicture}
